\documentclass[11pt]{article}
\usepackage[margin=1in]{geometry}
\usepackage[T1]{fontenc}
\usepackage{lmodern,microtype}
\usepackage{amsmath,amssymb,amsthm,mathtools}
\usepackage{booktabs,array,graphicx,xcolor,url}
\usepackage[colorlinks=true,linkcolor=blue!55!black,citecolor=blue!55!black,urlcolor=blue!55!black]{hyperref}

\renewcommand{\topfraction}{.9}\renewcommand{\bottomfraction}{.6}
\renewcommand{\textfraction}{.08}\renewcommand{\floatpagefraction}{.8}

\newtheorem{theorem}{Theorem}
\newtheorem{proposition}[theorem]{Proposition}
\newtheorem{lemma}[theorem]{Lemma}
\newtheorem{corollary}[theorem]{Corollary}
\theoremstyle{remark}
\newtheorem{remark}[theorem]{Remark}

\newcommand{\one}{\mathbf 1}
\newcommand{\E}{\mathbb E}
\newcommand{\PP}{\mathbb P}
\newcommand{\N}{\mathbb N}
\newcommand{\diag}{\operatorname{diag}}
\newcommand{\Cov}{\operatorname{Cov}}
\newcommand{\Var}{\operatorname{Var}}
\newcommand{\TV}{\mathrm{TV}}
\newcommand{\Hm}{\mathcal H}
\newcommand{\gen}{\mathcal L}
\newcommand{\Bin}{\operatorname{Bin}}

\title{Conservative molecular inheritance and population risk:\\
extinction and variance ordering, approximation limits, and regrowth}
\author{}
\date{}

\begin{document}
\maketitle

\begin{abstract}
When a dividing cell distributes a conserved set of molecular marks between
its daughters, the two sister states are dependent even though their later
lineages evolve independently. We study a finite reader--writer chromatin
motif with $N$ modifiable sites, explicit cell division and a
state-dependent death hazard, and compare the exact complementary
(conservative) allocation of old marks with the common simplification that
draws the two daughters independently from the same marginal law. The two
models have identical type-resolved mean dynamics. We prove that, for every
$N$, every choice of nonnegative rate constants, constant division rate and
any death hazard that is monotone in the natural order on mark states,
complementary allocation gives the smaller extinction probability, the
smaller population generating function at every time, and the smaller
population variance; under irreducibility the extinction inequality is strict
at every founder state whenever survival is possible. An exact identity
writes the extinction difference as the conditional covariance of daughter
continuation risks propagated by a nonnegative response operator, and
positive-vector bounds make it quantitative. The raw pair laws are far apart
in total variation, so any accuracy of the independence approximation is
observable-specific. Exact rational certificates and a validated
Taylor integration give, for a two-site source, a difference of more than
$0.098$ in the probability of ever reaching $300$ cells and of more than
$0.084$ in reaching $300$ cells by time $300$. For a $16$-site source with a
sigmoidal hazard, a near-balanced founder has an extinction gap above
$0.0128$ uniformly over the slope interval $[7,9]$, while the all-active
founder has a gap below $3.6\times10^{-4}$. For a fixed smooth readout we
prove that survival from balanced founders tends to zero as $N\to\infty$ at
an explicit polynomial rate, so the certified finite-size contrast is a
finite-size effect; a discontinuous readout behaves differently and is left
open, as is state-dependent division, for which we exhibit an exact
obstruction to the proof. We also derive what paired-clone and barcode
observations can and cannot identify. Probability arguments are conventional;
selected finite algebra is checked in Lean, and all numerical claims are exact
rational comparisons. The examples are mechanistic and are not clinically
calibrated.
\end{abstract}

\section{Introduction}

A mean growth law cannot say whether a small population dies out or reaches a
consequential size. That question depends on the joint law of the offspring,
and in particular on how a mother's state is divided between her daughters.
Non-genetic cellular memory is often carried by a finite stock of molecular
material: modified nucleosomes are recycled to the two daughter strands at
the replication fork in roughly equal shares
\cite{alabert2012,petryk2018,yu2018}, and random partitioning of molecules at
division is a major source of phenotypic heterogeneity
\cite{huh2011ng,huh2011pnas}. Chromatin-mediated drug-tolerant states
\cite{sharma2010,shaffer2017,oren2021} make the link to population survival
concrete: a lineage persists only while enough of its cells remain in a
favourable molecular state.

Conservation has a simple probabilistic consequence. If one daughter inherits
many old activating marks, her sister inherits few. The two sister
\emph{states} are therefore negatively dependent, even in a model in which
their future subtrees are conditionally independent. A modeller who knows
only the law of a single daughter may be tempted to draw both daughters
independently from that marginal. This \emph{independent comparator}
preserves every first moment of the population, but it generally violates
conservation of the mother's material across the pair. We ask:
\begin{quote}
\emph{Does ignoring the sister dependence change population risk in a
definite direction? How large is the change, and which preparations, readouts
and clocks govern it?}
\end{quote}

\paragraph{Results.}
We work with a finite reader--writer motif with recruited erasure
\cite{dodd2007,bruno2022,bruno2024} on $N$ sites, a constant division rate
$b$, and a death hazard $d$ that decreases with activating and increases with
repressive marks. Section~\ref{sec:source} constructs the chronological
branching population and its embedded offspring maps.
\begin{enumerate}
\item \emph{Direction} (Theorem~\ref{thm:direction}). For every $N$ and every
founder state, complementary allocation has the smaller extinction
probability and the smaller scalar population generating function at every
time. The same holds for finitely many molecular-independent Erlang division
phases. The proof combines a monotone coupling of the molecular chain, an
order-preserving coupling of allocations, and Harris' inequality
\cite{harris1960,esary1967} after reversing the repressive allocation bits.
\item \emph{Variance} (Corollary~\ref{cor:variance}). The population size has
the same mean and a smaller variance under complementary allocation, at every
time; the second factorial moment obeys a linear equation forced by the
covariance of daughter continuation \emph{means}.
\item \emph{Strictness} (Proposition~\ref{prop:strict}). With positive basal
rates and a nonconstant hazard there is a dichotomy: either both models die
out almost surely, or the extinction inequality is strict at every state.
\item \emph{Magnitude} (Theorem~\ref{thm:response},
Proposition~\ref{prop:limits}). The extinction gap equals a nonnegative
response operator applied to the conditional covariance of daughter
continuation \emph{risks}. Weighted contraction bounds enclose it. The pair
laws themselves are at total-variation distance
$1-\binom{2a}{a}\binom{2r}{r}4^{-(a+r)}$, which tends to one, so closeness of
the two models can only be observable-specific.
\item \emph{Certified instances} (Section~\ref{sec:instances}). Exact rational
risk boxes give $0.1144144<q_I-q_J<0.1144331$ for an all-active two-site
founder, and a $16$-site sigmoidal-hazard certificate that is uniform over
the slope interval $[7,9]$.
\item \emph{Large $N$} (Theorem~\ref{thm:boundary}). For a fixed Lipschitz
readout that is subcritical on the symmetry line, survival from founders
within $O(N^{-1/2})$ of that line is at most $O(N^{-\gamma}(\log N)^{3/2})$
with an explicit $\gamma>0$, in both models. The proof is a first-moment
(many-to-one) argument and uses no population coupling.
\item \emph{Regrowth} (Theorems~\ref{thm:threshold}--\ref{thm:application}).
An exact identity converts extinction boxes into bounds on reaching large
population thresholds, and a validated integration of the microscopic
generating function gives a finite-deadline statement.
\item \emph{Observation} (Section~\ref{sec:observation}). Paired sister-clone
extinction has nonpositive conditional covariance; we give a sensitivity
interval under parental mixing, misclassification and missingness, a
finite-sample radius, and the exact effect of independent cell capture on
barcode-type clone counts.
\end{enumerate}

\paragraph{Relation to prior work.}
General multitype branching processes already allow an arbitrary joint law of
offspring types \cite{athreya1972,kimmel2015}; conditional independence of
subtrees does not require independent sister types. Crump and Mode studied
correlated sister lifetimes \cite{crump1969}, Olofsson developed branching
with local dependencies and showed that growth is governed by the marginals
\cite{olofsson1996}, and Duffy and Subramanian showed for lymphocyte models
that related-cell correlation leaves the mean unchanged while changing
variability \cite{duffy2009}; correlated division times are observed
experimentally \cite{sandler2015}. Sibling dependence through a shared random
environment \cite{vatutin2020} is a different mechanism. Hence ``correlation
changes risk but not the mean'' is not by itself new. Closest in spirit is the
division-within-division literature initiated by Kimmel \cite{kimmel1997},
where a cell's content is shared between daughters
\cite{bansaye2008,bansaye2011,marguet2024}; there the partitioning kernel
influences the fate of the content and of the cell population. Our
contribution is different: a \emph{comparison} theorem, for an explicit
chemical source, between the conservative joint law and its
independent-marginal replacement, with a sign, an exact magnitude identity,
certified finite instances and an asymptotic limit. On the chromatin side,
transfer-matrix treatments of partition and recovery
\cite{micheelsen2010} and replication-driven switching \cite{movilla2025}
are established; we do not claim a new switching asymptotic, and we keep
molecular reactions, death and division active on one clock rather than
assuming recovery between divisions. Models that connect inherited phenotype
to treatment response prescribe the inheritance rule \cite{cassidy2021}; here
it is derived from material conservation.

\paragraph{Evidence classes.}
Stochastic constructions, couplings and limit arguments are conventional
proofs given in full. Every printed numerical inequality is an exact
rational comparison produced by the accompanying scripts. A subset of the
finite algebra is additionally checked in Lean~4 with mathlib
\cite{moura2021,mathlib2020}; Section~\ref{sec:scope} states precisely which.
Floating-point curves are labelled as diagnostics.

\section{Molecular source and chronological population}\label{sec:source}

\subsection{Reader--writer source}
Fix $N\ge1$ and let $S_N=\{(a,r)\in\N^2:a+r\le N\}$, with $u=N-a-r$.
A state records the numbers of activating ($A$), repressive ($R$) and
unmodified ($U$) sites. The molecular generator $Q$ has the four
nearest-neighbour rates
\begin{align}
 q_{(a,r),(a+1,r)}&=u\,(w_A+k_Aa/N),&q_{(a,r),(a,r+1)}&=u\,(w_R+k_Rr/N),\label{eq:write}\\
 q_{(a,r),(a-1,r)}&=a\,(e_A+h_Ar/N),&q_{(a,r),(a,r-1)}&=r\,(e_R+h_Ra/N),\label{eq:erase}
\end{align}
with nonnegative constants: basal writing $w$, read--write feedback $k$,
intrinsic erasure $e$ and erasure recruited by the opposite mark $h$. Vanishing
counts enforce the boundary of the triangle; diagonal entries of $Q$ are minus
the outgoing rates; $|S_N|=(N+1)(N+2)/2$. The factors $1/N$ fix how effective
interactions scale with the number of sites; a different contact scaling is a
different model. The names ``activating'' and ``repressive'' designate a
modelled functional order and assert nothing universal about particular
chromatin marks. The constants $e_A,e_R$ are \emph{intrinsic} erasure rates:
dilution by replication is represented explicitly by division below, so an
effective erasure rate from an undivided model that already averages
replication loss \cite{bruno2022} must not be imported, or that loss would be
counted twice.

\subsection{Complementary and independent daughters}
At division each old mark independently goes to daughter one with probability
$1/2$, and new unmodified sites restore each daughter to $N$ sites. For a
mother $x=(a,r)$ the daughter pair is
\begin{equation}\label{eq:partition}
 (Y,Z)=\bigl((I,J),(a-I,r-J)\bigr),\qquad I\sim\Bin(a,\tfrac12),\ J\sim\Bin(r,\tfrac12)\ \text{independent}.
\end{equation}
Write $D_1f(x)=\E_xf(Y)=\E_xf(Z)$ and $D_J(f,g)(x)=\E_x[f(Y)g(Z)]$. The
\emph{independent comparator} draws two independent copies of $Y$, so that
$D_I(f,g)=(D_1f)(D_1g)$. It has the same marginals, but it does not conserve
the mother's marks across the pair: it is an approximation to be evaluated,
not an alternative physical redistribution. The vector $(I,a-I)$ is a
two-cell multinomial and is negatively associated in the sense of
\cite{joagdev1983}; what we need below is a sign for functions monotone in a
\emph{mixed} order, which requires the argument of Lemma~\ref{lem:assoc}.

\subsection{Population process}
A cell of type $x$ evolves by $Q$, divides at constant rate $b>0$ and dies at
rate $d(x)\ge0$. At division it is replaced by two daughters with joint law
$D\in\{D_J,D_I\}$; the subtrees of distinct cells are independent given their
birth types. This prescribes a continuous-time Markov chain on finite
configurations of typed cells.

\begin{lemma}[Construction]\label{lem:construction}
For either daughter law and any finite founder configuration:
\begin{enumerate}
\item[(i)] The process is nonexplosive, its size $K(t)$ is stochastically
dominated by a Yule process of rate $b$, and $\sup_{s\le t}K(s)$ has finite
moments of all orders.
\item[(ii)] With
\begin{equation}\label{eq:R}
 \Hm=\diag(b+d)-Q,\qquad R=\Hm^{-1}=\int_0^\infty e^{(Q-\diag(b+d))t}\,dt\ \ge0,\qquad T=bR,
\end{equation}
the extinction probabilities $q=(q_x)_{x\in S_N}$ from one founder are the
least fixed point in $[0,1]^{S_N}$ of the offspring map
\begin{equation}\label{eq:F}
 F_J(v)=Rd+T\,D_J(v,v)\qquad\text{or}\qquad F_I(v)=Rd+T\,(D_1v)^2 ,
\end{equation}
and $q=\lim_nF^n(0)$. Squares of vectors are coordinatewise. Any
$\bar v\in[0,1]^{S_N}$ with $F(\bar v)\le\bar v$, equivalently
$\Hm\bar v\ge d+bD(\bar v,\bar v)$, satisfies $q\le\bar v$.
\item[(iii)] For $0\le z\le1$ the generating function
$u_x(t;z)=\E_xz^{K(t)}$ \textup{(}$0^0=1$\textup{)} is the unique solution of
\begin{equation}\label{eq:ODE}
 \dot u=Qu+d\,(\one-u)+b\,\{D(u,u)-u\},\qquad u(0;z)=z\one ,
\end{equation}
which leaves $[0,1]^{S_N}$ invariant; $u_x(t;0)$ is the probability of
extinction by time $t$.
\item[(iv)] $m_x(t)=\E_xK(t)$ satisfies $\dot m=Am$, $m(0)=\one$, with the
\emph{common} mean generator
\begin{equation}\label{eq:mean}
 A=Q+b\,(2D_1-I)-\diag(d).
\end{equation}
The same holds for every type-resolved first moment, so the two models have
identical mean dynamics.
\end{enumerate}
\end{lemma}
\begin{proof}
(i) Each cell divides at rate $b$ whatever its type and deaths only remove
cells, so a coupling with common division clocks bounds $K$ by a Yule
process, whose value at time $t$ is geometric. On $\{K\le n\}$ all rates are
bounded; localisation at the hitting time of size $n$ and the Yule bound give
nonexplosion and the moment bound.

(ii) The killed semigroup $e^{(Q-\diag(b+d))t}$ is nonnegative with row sums
at most $e^{-bt}$, which gives \eqref{eq:R} and $R\ge0$. Starting from $x$,
$R_{xy}d_y$ and $R_{xy}b$ are the probabilities that the first terminal event
of the founder is a death, respectively a division, occurring in molecular
state $y$. Conditioning on it and using independence of the two subtrees
given $(Y,Z)$ shows that $F^n(0)_x$ is the probability that the family tree
of the founder has fewer than $n$ generations. Since $\Hm\one=b\one+d$ we have
$F(\one)=\one$: death mass is retained. The monotone limit $q^*$ is the least
fixed point of the polynomial map $F$ and lies below every supersolution.
Every cell has a terminal event almost surely because $b>0$, so a finite tree
dies out in finite time; conversely an infinite, locally finite binary tree
contains an infinite line of descent, which by nonexplosion is not completed
in finite time. Thus $q^*$ is the probability of eventual extinction and not
merely an algebraic root. The equivalence of the two supersolution conditions
follows on multiplying by $R\ge0$.

(iii) Condition on the first event in $[0,\delta]$; by (i) two or more events
have probability $O(\delta^2)$ uniformly in the founder type, and after a
division the branching property gives $u_Y(t;z)u_Z(t;z)$ averaged over the
\emph{joint} law of $(Y,Z)$. The vector field is polynomial, hence locally
Lipschitz, and points inward on the faces of the cube. Zero is absorbing, so
$u_x(t;0)=\PP_x(K(t)=0)$.

(iv) Differentiate \eqref{eq:ODE} at $z=1$, which is justified by (i), and
use $D(f,\one)=D(\one,f)=D_1f$ for both laws. For a nonnegative weight
$\varphi$ on types the same computation applied to
$\E_x\sum_{\text{cells}}\varphi$ gives $e^{tA}\varphi$.
\end{proof}

\section{Direction of the inheritance effect}\label{sec:direction}

Order the states by
\[
 (a,r)\preceq(a',r')\quad\Longleftrightarrow\quad a\le a'\ \text{and}\ r\ge r' .
\]
An upper state has more activating and less repressive material. A function
on $S_N$ is called decreasing (increasing) if it is nonincreasing
(nondecreasing) in $\preceq$.

\begin{lemma}[Monotone molecular coupling]\label{lem:molecular}
For any nonnegative constants in \eqref{eq:write}--\eqref{eq:erase} and any
$x\preceq x'$ there is a Markovian coupling $(X_t,X'_t)$ of two copies of the
chain $Q$ started at $(x,x')$ with $X_t\preceq X'_t$ for all $t$.
\end{lemma}
\begin{proof}
Couple each of the four channels at the minimum of its two rates, letting the
excess rate fire in one chain only. A joint jump leaves both coordinate
differences unchanged. A solitary jump changes one coordinate difference by
one, so it can break the order only if that coordinate is currently equal in
the two chains. Let $x=(a,r)\preceq x'=(a',r')$. If $a=a'$ then $u\le u'$
because $r\ge r'$, so the activating-write rate $u(w_A+k_Aa/N)$ of the lower
chain does not exceed that of the upper chain, and the activating-erasure
rate $a(e_A+h_Ar'/N)$ of the upper chain does not exceed that of the lower
chain. If $r=r'$ then $u'\le u$, so the repressive-write rate of the upper
chain does not exceed that of the lower one, and the repressive-erasure rate
$r(e_R+h_Ra/N)$ of the lower chain does not exceed that of the upper one.
In each of the four cases the excess jump therefore occurs in the chain for
which it preserves the order. At $a+r=N$ the write rates vanish, and at $a=0$
or $r=0$ the corresponding erasure rate vanishes, so the coupling respects the
triangle.
\end{proof}

\begin{lemma}[Ordered allocation and population embedding]\label{lem:embedding}
Assume $b$ is constant and $d$ is decreasing. For $x\preceq x'$ and either
daughter law, the populations started from one founder of type $x$ and one of
type $x'$ can be coupled so that at every time there is an injection from the
living cells of the first into the living cells of the second under which
types are ordered by $\preceq$. Consequently, in each model,
$x\mapsto u_x(t;z)$ and $x\mapsto q_x$ are decreasing, $x\mapsto m_x(t)$ is
increasing, and $F_J$, $F_I$ map the cone
\[
 \mathcal C=\{f:S_N\to[0,1]\ \text{decreasing}\}
\]
into itself. Both maps are also monotone for the coordinatewise order.
\end{lemma}
\begin{proof}
Run a matched pair of cells with the coupling of Lemma~\ref{lem:molecular},
a common division clock of rate $b$, a common death clock at the smaller
hazard $d(X'_t)$ and an additional death clock of rate $d(X_t)-d(X'_t)\ge0$
for the lower cell alone. If the lower cell dies alone, the upper cell and
its descendants continue unmatched. At a common division with mother states
$(a,r)\preceq(a',r')$ use shared fair bits for the $a$ common activating
marks and the $r'$ common repressive marks, and fresh independent bits for
the $a'-a$ additional activating marks of the upper mother and the $r-r'$
additional repressive marks of the lower mother. If $(I,J')$ counts the
shared marks sent to daughter one and $I'',J''$ the additional ones, the
daughters are
\[
 (I,\,J'+J'')\preceq(I+I'',\,J')\qquad\text{and}\qquad
 (a-I,\,r-J'-J'')\preceq(a'-I-I'',\,r'-J') ,
\]
the second relation because $I''\le a'-a$ and $J''\le r-r'$. Both daughter
pairs are thus ordered and each mother sees her own complementary law. For
independent daughters apply the same marginal coupling separately and
independently to daughter one and to daughter two. Recursion over matched
pairs, with nonexplosion from Lemma~\ref{lem:construction}, gives the
embedding; hence $K^{x}(t)\le K^{x'}(t)$ pathwise, which yields the three
monotonicity statements. For the cone, stop the matched founders at the first
terminal event and pay $1$ for a death and $f(Y)f(Z)$ for a division: a
solitary death of the lower cell pays $1\ge$ any payoff of the upper cell, a
common death pays $1$ to both, and a common division pays
$f(Y)f(Z)\ge f(Y')f(Z')$ because $f$ is decreasing and nonnegative.
Coordinatewise monotonicity holds because $R,T$ and the coefficients of $D$
are nonnegative.
\end{proof}

\begin{lemma}[Sign of the sister covariance]\label{lem:assoc}
If $f:S_N\to\mathbb R$ is monotone in $\preceq$ \textup{(}in either
direction\textup{)}, then for every mother state $x$
\begin{equation}\label{eq:covsign}
 \Cov_x\bigl(f(Y),f(Z)\bigr)\le0,\qquad\text{that is,}\qquad D_J(f,f)\le(D_1f)^2 .
\end{equation}
\end{lemma}
\begin{proof}
Let $\alpha_1,\dots,\alpha_a$ and $\beta_1,\dots,\beta_r$ be the independent
fair bits indicating that an activating, respectively repressive, mark goes
to daughter one, and put $\tilde\beta_j=1-\beta_j$. Then
$Y=(\sum\alpha_i,\;r-\sum\tilde\beta_j)$ and
$Z=(a-\sum\alpha_i,\;\sum\tilde\beta_j)$, so in the order $\preceq$ the state
$Y$ is increasing and $Z$ is decreasing in each of the independent bits
$(\alpha,\tilde\beta)$. Hence $f(Y)$ and $f(Z)$ are monotone in opposite
directions in every bit. Harris' inequality \cite{harris1960,esary1967}
states that two functions of independent variables that are increasing in
each variable have nonnegative covariance; apply it to $-f(Y)$ and $f(Z)$
or to $f(Y)$ and $-f(Z)$. For completeness: condition on the last bit
$\xi$. For functions $G,G'$ that are increasing in each bit, the law of total
covariance gives
$\Cov(G,G')=\E\Cov(G,G'\mid\xi)+\tfrac14(g_1-g_0)(g'_1-g'_0)$, where $g_i,g'_i$
are the conditional means given $\xi=i$; the first term is nonnegative by
induction on the number of bits and the second because both factors are
nonnegative.
\end{proof}

The cone is essential: for an all-active mother with two marks,
$f(0)=f(2)=1$, $f(1)=0$ gives covariance $+1/4$.

\begin{theorem}[Source comparison]\label{thm:direction}
Let $Q$ be given by \eqref{eq:write}--\eqref{eq:erase} with arbitrary
nonnegative constants, let $b>0$ be constant and let $d\ge0$ be decreasing in
$\preceq$. Then for every $N$, founder state $x$, $t\ge0$ and $z\in[0,1]$,
\[
 q_J(x)\le q_I(x),\qquad u_{J,x}(t;z)\le u_{I,x}(t;z).
\]
In particular $\PP_x(K_J(t)=0)\le\PP_x(K_I(t)=0)$ for all $t$. The same
conclusions hold when the exponential division clock is replaced by $m$
successive exponential phases whose rates do not depend on the molecular
state, with molecular reactions and mortality active in every phase, division
at the end of the last phase, and daughters born in the first phase.
\end{theorem}
\begin{proof}
\emph{Extinction.} Let $f_n=F_J^n(0)$ and $g_n=F_I^n(0)$. By
Lemma~\ref{lem:embedding}, $f_n\in\mathcal C$. If $f_n\le g_n$ then, by
\eqref{eq:covsign} and \eqref{eq:F} and then by monotonicity of $F_I$,
\[
 f_{n+1}=F_J(f_n)\le F_I(f_n)\le F_I(g_n)=g_{n+1}.
\]
Let $n\to\infty$ and use Lemma~\ref{lem:construction}(ii).

\emph{Generating functions.} Write \eqref{eq:ODE} as $\dot u=G_D(u)$. By
Lemma~\ref{lem:embedding}, $u_J(t;z)\in\mathcal C$ for all $t$, so
\eqref{eq:covsign} gives
$\dot u_J=G_J(u_J)\le G_I(u_J)$: the function $u_J$ is a subsolution of the
independent system with the same initial value. On the cube the
off-diagonal partial derivatives of $G_I$ are
$Q_{xy}+2b(D_1v)_xD_1(x,y)\ge0$, so $G_I$ is quasimonotone and the
Kamke--M\"uller comparison theorem \cite[Ch.~3]{smith1995} yields
$u_J\le u_I$.

\emph{Erlang clocks.} Append the phase $j\in\{1,\dots,m\}$ to the type and
compare types with equal phase only. Phase clocks do not depend on $(a,r)$
and are coupled identically, death is coupled as before in every phase, and
both daughters restart in phase one, so Lemmas~\ref{lem:molecular} and
\ref{lem:embedding} hold on the finite augmented space with the cone of
functions decreasing in $(a,r)$ at each phase. Lemma~\ref{lem:assoc} is
applied to $f(\cdot,1)$, and the two steps above are unchanged.
\end{proof}

This is an ordering of extinction probabilities and of scalar generating
functions. It is not stochastic dominance of population sizes, which for
nondegenerate cases would contradict the equality of means, and it does not
order every first-passage probability.

\subsection{Variance ordering}
\begin{corollary}[Equal means, smaller variance]\label{cor:variance}
Under the hypotheses of Theorem~\ref{thm:direction} the second factorial
moment $v_{D,x}(t)=\E_x[K(t)(K(t)-1)]$ is the solution of
\begin{equation}\label{eq:second}
 \dot v_D=Av_D+2b\,D(m,m),\qquad v_D(0)=0,
\end{equation}
where $m(t)=e^{tA}\one$ is the common mean. Consequently
\begin{equation}\label{eq:varorder}
 v_I(t)-v_J(t)=-2b\int_0^te^{(t-s)A}\bigl(\Cov_x(m_Y(s),m_Z(s))\bigr)_x\,ds\ \ge0,
 \qquad \Var_xK_J(t)\le\Var_xK_I(t)
\end{equation}
for every $x$ and $t$. The inequality persists for any common initial mixture
of founder types and for sums over finitely many independent founders.
\end{corollary}
\begin{proof}
By Lemma~\ref{lem:construction}(i) the generating function is twice
differentiable at $z=1$. Differentiating \eqref{eq:ODE} twice at $z=1$, where
$u=\one$, $\partial_zu=m$, and using
$\partial_z^2D(u,u)=2D_1\partial_z^2u+2D(m,m)$, gives \eqref{eq:second}. By
Lemma~\ref{lem:embedding} $m(s)$ is increasing, so Lemma~\ref{lem:assoc}
gives $(D_1m)^2-D_J(m,m)\ge0$; the matrix $A$ has nonnegative off-diagonal
entries, so $e^{(t-s)A}\ge0$ and variation of constants gives
\eqref{eq:varorder}. The means agree, so the variances are ordered. For a
mixture the between-founder variance of the conditional mean is common to the
two models, and for independent founders variances add.
\end{proof}
Alternatively, $u_I(t;z)-u_J(t;z)\ge0$ on $[0,1]$ vanishes together with its
first derivative at $z=1$, which forces a nonnegative second derivative
there. The mechanism is visible in \eqref{eq:varorder}: the covariance of
daughter continuation \emph{means} forces the variance, whereas
Section~\ref{sec:response} shows that the covariance of daughter continuation
\emph{risks} forces the extinction gap. Smaller variance at equal mean does
not in general imply smaller extinction probability for branching processes;
here both follow from the same association argument, and neither orders
higher moments or first-passage events. Figure~\ref{fig:summary}C shows the
variance ratio for the two worked sources of Section~\ref{sec:instances}: at
$t=60$ it is $91.9/127.6\approx0.72$ for the two-site source and
$15510/16276\approx0.95$ for the $16$-site source (diagnostic values).

\subsection{Strictness}
The inequalities of Theorem~\ref{thm:direction} are not strict in general:
if both models die out almost surely then $q_J=q_I=\one$. Under mild
conditions this is the only exception.

\begin{proposition}[Dichotomy]\label{prop:strict}
In addition to the hypotheses of Theorem~\ref{thm:direction}, assume
$w_A,w_R,e_A,e_R>0$ and that $d$ is not constant. Let $M=2TD_1$ be the common
mean offspring matrix of the embedded generation process and $\rho(M)$ its
spectral radius. If $\rho(M)\le1$ then $q_J=q_I=\one$. If $\rho(M)>1$ then
\[
 q_J(x)<q_I(x)<1\qquad\text{for every }x\in S_N .
\]
\end{proposition}
\begin{proof}
Positive basal rates make $Q$ irreducible on $S_N$, so $R$ and $T$ are
entrywise positive, and $D_1(y,y)=2^{-(a+r)}>0$ makes $M$ entrywise positive.
The generation process is a multitype Galton--Watson process, with joint
offspring types under $D_J$ and independent ones under $D_I$, which is
positively regular and nonsingular with mean matrix $M$ in both models. By the
classical criterion \cite[Ch.~V.3]{athreya1972}, $q=\one$ if $\rho(M)\le1$
and $q<\one$ in every coordinate if $\rho(M)>1$.

Assume $\rho(M)>1$. If $q_J$ were a constant $c<1$, then
$\Hm q_J=d+bD_J(q_J,q_J)$ would read $(b+d_x)c=d_x+bc^2$, that is
$(1-c)(d_x-bc)=0$, for every $x$, contradicting that $d$ is not constant. So
the decreasing function $q_J$ is not constant, and since the path
$(0,N),(0,N-1),\dots,(0,0),(1,0),\dots,(N,0)$ is increasing in $\preceq$,
the function $q_J$ is nonconstant on $\{(i,0)\}_{i\le N}$ or on
$\{(0,j)\}_{j\le N}$. In the first case take the mother $(N,0)$. The last
step of the proof of Lemma~\ref{lem:assoc} applied to $G=-q_J(Y)$,
$G'=q_J(Z)$ gives
\[
 \Cov_{(N,0)}\bigl(q_J(Y),q_J(Z)\bigr)\le-\tfrac14\Bigl(\E\bigl[q_J(I',0)-q_J(I'+1,0)\bigr]\Bigr)^2<0,
 \qquad I'\sim\Bin(N-1,\tfrac12),
\]
where the two conditional-mean differences coincide by the symmetry
$I'\mapsto N-1-I'$, and the expectation is strictly positive because $I'$
charges every value in $\{0,\dots,N-1\}$ and $q_J(\cdot,0)$ is decreasing and
nonconstant. The second case is identical with the mother $(0,N)$. By
\eqref{eq:forcing} below the forcing
$\delta=-T(\Cov_x(q_J(Y),q_J(Z)))_x$ is then strictly positive at
\emph{every} state because $T>0$, and
Theorem~\ref{thm:response} gives $q_I-q_J\ge\delta>0$.
\end{proof}
If $d$ is constant the type is irrelevant and $q_J=q_I=\min(1,d/b)\one$.

\subsection{State-dependent division: an obstruction and an open question}
\label{sec:statedep}
For a bounded state-dependent division rate $b(x)$ the construction of
Section~\ref{sec:source} goes through with $T=R\diag(b)$ and
$A=Q+\diag(b)(2D_1-I)-\diag(d)$, and the two models still share their mean.
The direction theorem, however, rests on cone preservation, and
monotonicity of $b$ does not restore it: a faster division clock in a
well-marked cell spreads its marks over two less-marked cells sooner.

\begin{proposition}[Exact cone obstruction]\label{prop:obstruction}
Let $N=4$, $Q=0$, $d$ constant, and for some $c\in(0,1)$ let $f(a,r)=1$ for
$a\le2$ and $f(a,r)=c$ for $a\ge3$, a decreasing function. Then
\[
 D_J(f,f)(3,0)-c=\tfrac34(1-c),\qquad D_J(f,f)(4,0)-c=\tfrac38(1-c),
\]
so if $b$ is increasing with $b(4,0)>2\,b(3,0)$, the vector field of
\eqref{eq:ODE} at $f$ is strictly larger at $(4,0)$ than at $(3,0)$ although
$f(3,0)=f(4,0)$ and $(3,0)\preceq(4,0)$: it points out of $\mathcal C$.
\end{proposition}
\begin{proof}
For the mother $(3,0)$ the product $f(Y)f(Z)$ equals $c$ exactly when one
daughter receives all three marks, which has probability $1/4$, and equals
$1$ otherwise. For $(4,0)$ it equals $c$ unless the split is $(2,2)$, which
has probability $3/8$. The molecular and death terms of the vector field
coincide at the two states.
\end{proof}
This is a failure of the proof route at a test vector; it is not an example
with $q_J>q_I$. A random search followed by hill climbing over roughly
$8\times10^3$ instances with $N\le5$, $b$ increasing and $d$ decreasing,
including $Q=0$, found no reversal of the extinction order beyond floating
point noise (diagnostic). Whether Theorem~\ref{thm:direction} survives
state-dependent division under a natural condition, perhaps only along the
actual generating-function flow, is open. The Erlang extension of
Theorem~\ref{thm:direction} accommodates a variable cell-cycle duration, but
not a trade-off between proliferation and molecular state.

\section{Magnitude: covariance and population response}\label{sec:response}

For any vector $v$, conditioning on the complete pre-division mother state
gives
\begin{equation}\label{eq:forcing}
 F_J(v)-F_I(v)=T\bigl(\Cov_x(v(Y),v(Z))\bigr)_{x\in S_N}.
\end{equation}
The resolvent $T$ weights the molecular states actually visited before
division, and the covariance weights daughters by their \emph{continuation
risk}, not by raw mark abundance.

\begin{theorem}[Exact response and enclosures]\label{thm:response}
Let $g=q_I-q_J$,
\[
 \delta=F_I(q_J)-F_J(q_J)=-T\bigl(\Cov_x(q_J(Y),q_J(Z))\bigr)_x,\qquad
 B=T\diag\bigl(D_1(q_I+q_J)\bigr)D_1 .
\]
Then $g=\delta+Bg$, and under Theorem~\ref{thm:direction} $\delta\ge0$, so
$g\ge\sum_{k=0}^nB^k\delta$ for every $n$; if $\rho(B)<1$,
\begin{equation}\label{eq:response}
 g=(I-B)^{-1}\delta .
\end{equation}
If $q_I\le\bar q$ and a vector $\nu>0$ satisfies
$DF_I(\bar q)\nu\le\kappa\nu$ with $\kappa<1$, where
$DF_I(v)=2T\diag(D_1v)D_1$, then
\begin{equation}\label{eq:margin}
 \|g\|_\nu\le\frac{\|\delta\|_\nu}{1-\kappa},\qquad \|v\|_\nu=\max_x|v_x|/\nu_x .
\end{equation}
Given boxes $0\le\ell_J\le q_J\le v_J\le\one$ and
$0\le\ell_I\le q_I\le v_I\le\one$, put
\begin{align*}
 \delta_-&=\max\{0,\,T[(D_1\ell_J)^2-D_J(v_J,v_J)]\},&
 \delta_+&=T[(D_1v_J)^2-D_J(\ell_J,\ell_J)],\\
 B_-&=T\diag(D_1(\ell_I+\ell_J))D_1,&B_+&=T\diag(D_1(v_I+v_J))D_1 .
\end{align*}
If $B_+\nu\le\kappa_+\nu$ for some $\nu>0$, $\kappa_+<1$, then
\begin{equation}\label{eq:box}
 (I-B_-)^{-1}\delta_-\ \le\ g\ \le\ (I-B_+)^{-1}\delta_+ .
\end{equation}
\end{theorem}
\begin{proof}
Subtract $q_J=F_J(q_J)$ from $q_I=F_I(q_I)$ and insert $\pm F_I(q_J)$:
$g=\delta+T[(D_1q_I)^2-(D_1q_J)^2]=\delta+Bg$ by the difference of two
squares. Iterating with $B\ge0$, $\delta\ge0$, $g\ge0$ gives the partial-sum
bound. Entrywise $0\le B\le DF_I(q_I)\le DF_I(\bar q)$, so the induced
$\nu$-norm of $B$ is at most $\kappa$, the Neumann series converges to a
nonnegative inverse of norm at most $(1-\kappa)^{-1}$, and
\eqref{eq:response}--\eqref{eq:margin} follow. The polynomial coefficients
are nonnegative, so $\delta_-\le\delta\le\delta_+$ and $B_-\le B\le B_+$
entrywise; termwise comparison of Neumann series gives \eqref{eq:box}.
\end{proof}

The factor $(1-\kappa)^{-1}$ is an explicit dependence on the distance from
criticality. It does not mean that the gap diverges there: as criticality is
approached the forcing $\delta$ decreases, and in the subcritical regime it
vanishes. We prove this margin-dependent bound and make no claim of a uniform
near-critical expansion. The premise concerns a Jacobian bound on a valid
\emph{enclosing box}; a Jacobian evaluated at an approximate numerical root
would not suffice.

\begin{proposition}[Limits of an independence approximation]\label{prop:limits}
For a mother $(a,r)$,
\begin{equation}\label{eq:tv}
 \|D_J-D_I\|_{\TV}=1-\binom{2a}{a}\binom{2r}{r}4^{-(a+r)} .
\end{equation}
If changing one allocation bit changes $f(Y)$ by at most $L/N$, then
\begin{equation}\label{eq:regularity}
 \bigl|\Cov_x(f(Y),f(Z))\bigr|\le\frac{(a+r)L^2}{4N^2}\le\frac{L^2}{4N}.
\end{equation}
If this holds for $f=q_J$ with $L=L_N$, and \eqref{eq:margin} holds, then
$\|q_I-q_J\|_\nu\le\|T\one\|_\nu\,L_N^2/[4N(1-\kappa)]$.
\end{proposition}
\begin{proof}
With $p_{ij}=\PP(I=i,J=j)$, the joint law puts mass $p_{ij}$ on the pair
$((i,j),(a-i,r-j))$ and nothing off this conservation graph, while the
product law puts mass $p_{ij}p_{a-i,r-j}=p_{ij}^2$ there. Hence the distance
is $1-\sum p_{ij}^2$, and $\sum_i\binom ai^2=\binom{2a}a$ gives
\eqref{eq:tv}. For \eqref{eq:regularity}, the Efron--Stein inequality
\cite{efron1981} for a function of $a+r$ independent fair bits with bounded
differences $L/N$ gives $\Var f(Y)\le(a+r)L^2/(4N^2)$; the same bound holds
for $f(Z)$ because $Z$ has the same marginal law, and Cauchy--Schwarz
concludes. The last assertion follows from $\delta\le T\one\cdot L_N^2/(4N)$
and \eqref{eq:margin}.
\end{proof}

For an all-active mother the distance \eqref{eq:tv} is
$1-\binom{2N}N4^{-N}\sim1-(\pi N)^{-1/2}\to1$: fractional concentration does
not make the pair laws close, and no argument through small total variation
can justify the independent comparator. What can be small is the covariance
of a \emph{regular} observable. A bounded regularity constant $L_N$, bounded
$\|T\one\|_\nu$ and a stability margin bounded away from zero would give an
$O(1/N)$ error. These are premises, not established uniform properties of
this source: for the smooth readout of Section~\ref{sec:instances} the
diagnostic constant $N\max|q_J(a+1,r)-q_J(a,r)|$ grows from $1.2$ at $N=2$ to
$9.0$ at $N=32$, because the extinction probability varies rapidly near the
symmetry line.

\section{Certified source instances}\label{sec:instances}

Throughout this section the constants are symmetric,
$w=0.01$, $k=1$, $h=0.5$, $e=0.01$, $b=0.1$ (Table~\ref{tab:param}). The
\emph{step} readout has $d=0.01$ when $a>r$ and $d=0.3$ otherwise. The
\emph{smooth} readout is
\begin{equation}\label{eq:smooth}
 d_{N,s}(a,r)=0.01+\frac{0.29}{1+\exp\bigl(s(a-r)/N\bigr)},\qquad s>0 .
\end{equation}
Both are decreasing in $\preceq$, so Theorem~\ref{thm:direction} applies to
every $N$ and every $s>0$. Smoothing changes intermediate hazards as well as
removing the discontinuity, and neither map is fitted to data.

\begin{table}[t]
\centering\small
\begin{tabular}{p{.13\linewidth}p{.20\linewidth}p{.57\linewidth}}
\toprule
Quantity & Value & Provenance and role\\\midrule
$N$ & $2$ or $16$ sites & Mathematical witnesses; not a measured locus size.\\
$w,e$ & $0.01$ per time & Symmetric basal writing and \emph{intrinsic} erasure; exploratory.\\
$k,h$ & $1,\ 0.5$ per time & Feedback at unit mark fraction and recruited erasure; contact factor $1/N$.\\
$b$ & $0.1$ per time & Exponential division clock; mean cycle $10$ in the absence of death.\\
$d$ & $0.01$ to $0.3$ per time & Assumed functional mortality; not a fitted drug response.\\
$s$ & $[7.999,8.001]$; $[7,9]$ & Dimensionless slope; robustness intervals of the certificates.\\
$m$ & $1,2,4$ phases & Same-mean Erlang clocks with phase rate $m/10$ (diagnostic).\\
$H,\ L$ & $300$; $300$, $2000$ cells & Deadline and thresholds of Section~\ref{sec:regrowth}.\\\bottomrule
\end{tabular}
\caption{Parameters of the worked examples. Time is a declared model unit;
the ratios $k/b=10$, $d_{\max}/b=3$, $d_{\min}/b=0.1$ are what matter, and
rescaling time cannot calibrate them. No entry is a measured range.}
\label{tab:param}
\end{table}

\subsection{Two sites: exact risk response}
Order the six states as $UU,UR,RR,AU,AR,AA$, that is
$(0,0),(0,1),(0,2),(1,0),(1,1),(2,0)$, and use the step readout. Three
certificates play distinct roles.

\emph{Coarse supersolutions.} The rational vectors
\begin{align*}
 \bar q_I&=(.948115,\ .983103,\ .987779,\ .378814,\ .775706,\ .335634),\\
 \bar q_J&=(.94,\ .982,\ .987,\ .27,\ .735,\ .235)
\end{align*}
satisfy $\Hm\bar q\ge d+bD(\bar q,\bar q)$ for their respective laws; the
smallest independent residual is $435828959/(4\times10^{13})$. By
Lemma~\ref{lem:construction}(ii), $q_I\le\bar q_I<0.988\,\one$ and
$q_J\le\bar q_J$. Eight iterations of $F_I$ from zero, each rounded down to a
multiple of $10^{-6}$, give $q_I(AA)\ge0.313291$, hence already
$q_I(AA)-q_J(AA)>0.078$.

\emph{Contraction weight.} On the whole box $[0,\bar q_I]$,
\[
 \nu=(2.148,\ 2.011,\ 1.997,\ 4.880,\ 2.882,\ 4.895)\quad\text{satisfies}\quad DF_I(\bar q_I)\,\nu\le0.796\,\nu
\]
by exact evaluation, so \eqref{eq:margin} holds with $\kappa=0.796$.

\emph{Tight response.} Two hundred and fifty downward-rounded iterations
(denominator $10^{12}$) give lower boxes; adding a small multiple of the
positive vector $(I-DF(\ell))^{-1}\one$ gives vectors whose supersolution
residuals are checked exactly. The resulting boxes have width below
$5\times10^{-6}$, for instance
$0.2207216<q_J(AA)<0.2207253$ and $0.3351447<q_I(AA)<0.3351497$, and
\eqref{eq:box} yields
\begin{equation}\label{eq:tightgap}
 0.1144144<q_I(AA)-q_J(AA)<0.1144331 .
\end{equation}
The lower response bound is strictly positive at all six states, from
$0.00246$ at $RR$ to $0.12347$ at $AU$, in agreement with
Proposition~\ref{prop:strict}. Neighbouring-state increments of $q_J$ are
bounded from the same boxes, so the regularity premise of
Proposition~\ref{prop:limits} is checkable at fixed $N$ without being
extrapolated.

\begin{figure}[tp]
\centering\includegraphics[width=.92\linewidth]{figures/fig_summary.pdf}
\caption{A: extinction gap against site count for all-active and balanced
founders under both readouts (floating-point diagnostics); the star is the
certified lower bound of Proposition~\ref{prop:smooth}(a) for the founder
$(9,7)$. B: survival of a balanced founder under complementary allocation;
under the smooth readout it decays, as Theorem~\ref{thm:boundary} requires
asymptotically, while under the step readout it does not (diagnostic; the
limit for the step readout is open). C: ratio of population variances,
which Corollary~\ref{cor:variance} bounds by one (diagnostic values).
D: proved bounds of Theorem~\ref{thm:application}, not estimates; the
independent upper bound for $300$ cells is an eventual bound and therefore
also bounds the deadline probability.}
\label{fig:summary}
\end{figure}

\begin{figure}[tbp]
\centering\includegraphics[width=.97\linewidth]{figures/fig_robustness.pdf}
\caption{A: certified lower bounds on $q_I(9,7)-q_J(9,7)$ on each of the $40$
slope cells of Proposition~\ref{prop:smooth}(b), with the floating-point gap
for comparison. B: extinction from $(9,7)$ at $N=16$, $s=8$ for Erlang
division clocks with the same mean (diagnostic); Theorem~\ref{thm:direction}
proves the sign for every $m$.}
\label{fig:robust}
\end{figure}

\subsection{Sixteen sites: readout slope and preparation}
\begin{proposition}[Smooth readout, $N=16$]\label{prop:smooth}
For $N=16$, the constants of Table~\ref{tab:param}, the exponential clock and
the hazard \eqref{eq:smooth}:
\begin{enumerate}
\item[(a)] for every $s\in[7.999,8.001]$,
\begin{gather*}
 0.48558095<q_J(9,7)<0.48561304,\qquad 0.49988773<q_I(9,7)<0.49991904,\\
 q_I(9,7)-q_J(9,7)>0.01427,\qquad 0\le q_I(16,0)-q_J(16,0)<0.000122,\\
 \max_xq_I(x)<0.997;
\end{gather*}
\item[(b)] for every $s\in[7,9]$,
\[
 q_I(9,7)-q_J(9,7)>0.0128,\qquad 0\le q_I(16,0)-q_J(16,0)<0.00036,\qquad \max_xq_I(x)<0.997 .
\]
\end{enumerate}
\end{proposition}
\begin{proof}
There are $153$ states. Fix a slope cell $[s_0,s_1]$. At a fixed state the
hazard \eqref{eq:smooth} is monotone in $s$, decreasing if $a>r$ and
increasing if $a<r$, so on the cell it lies between the smaller and the
larger of its two endpoint values. These endpoint values are enclosed by
bounding $e^y$, $0\le y\le9$, below by its degree-$120$ Taylor polynomial and
above by that polynomial plus a geometric bound on the tail, inverting for
negative arguments, and rounding outward at denominator $10^{14}$. This gives
rational vectors $d_-\le d_{N,s}\le d_+$ for all $s$ in the cell. Note that
$d_-$ and $d_+$ are statewise envelopes and neither is the hazard at an
endpoint slope: the hazards at two different slopes are not ordered, and
extinction need not be monotone in $s$.

A pointwise coupling of death clocks shows that in either model $q$ is
nondecreasing in the hazard vector. For $d_+$ we exhibit rational vectors
below one and verify every supersolution inequality
$\Hm\bar q\ge d_++bD(\bar q,\bar q)$ exactly. For $d_-$ we iterate from zero
the positive map
\begin{equation}\label{eq:G}
 G(z)_x=\frac{d_x+\sum_{y\ne x}Q_{xy}z_y+b\,D(z,z)_x}{d_x+b+\sum_{y\ne x}Q_{xy}},
\end{equation}
which has the same fixed points as \eqref{eq:F}, is monotone, and satisfies
$G(z)\le z'$ whenever $z\le z'$ and $z'$ is a fixed point; rounding all
nonnegative coefficients and iterates downward at denominator $10^{14}$ keeps
every iterate below the least fixed point by induction. For (a) a single
cell is used; $466$ joint and $474$ independent iterations give the displayed
bounds, and the largest box width over all $153$ coordinates and both models
is below $3.58\times10^{-5}$. For (b) the interval is divided into $40$ cells
of width $0.05$; the certified lower bound on the gap at $(9,7)$ increases
from $0.012847$ on $[7,7.05]$ to $0.014019$ on $[8.95,9]$
(Figure~\ref{fig:robust}A). All comparisons are exact rational comparisons,
and no solver tolerance is used.
\end{proof}

Enlarging the slope interval costs accuracy in a transparent way: a single
cell $[7,9]$ gives an envelope so wide that the resulting sufficient bound is
negative, although the true gap stays between $0.0138$ and $0.0147$
(diagnostic). A failed sufficient inequality is not evidence about the
dynamics; partitioning the interval repairs it.

The proposition is a source-specific obstruction to promoting an interior
estimate to a guarantee that is uniform over preparations: at the same $N$ an
all-active error below $1.3\times10^{-4}$ coexists with a near-balanced error
above $1.4\times10^{-2}$. Figure~\ref{fig:summary}A shows the corresponding
floating-point trends for $N=2,\dots,128$. From the all-active founder the
gap falls to $6\times10^{-10}$ (step) and $1.5\times10^{-6}$ (smooth) at
$N=128$. From the balanced founder $(N/2,N/2)$ it stays between $0.008$ and
$0.016$ for $8\le N\le128$. A third preparation, obtained by holding an
all-active cell for one time unit with division and death suspended, has law
$e^{Q^{\mathsf T}}\delta_{(N,0)}$ and gaps within one percent of the
all-active family for $N\le32$; this is a controlled pre-division hold, not a treatment
model. The certified preparation of Proposition~\ref{prop:smooth} is $(9,7)$,
the state of largest gap at $N=16$ under the smooth readout; it should not be confused with $(10,6)$
or $(8,8)$.


\subsection{Division-time variance}
For $m$ phases of rate $m/10$ let $V=(m/10)\bigl((m/10)I+\diag(d)-Q\bigr)^{-1}$;
successive phase conditioning gives the offspring map
$\one-V^m\one+V^mD(v,v)$. At $N=16$, $s=8$ and founder $(9,7)$, diagnostic
values of $(q_J,q_I)$ are $(0.48560,0.49990)$, $(0.51938,0.52860)$ and
$(0.53651,0.54262)$ for $m=1,2,4$, with gaps $0.01431$, $0.00922$, $0.00611$
and Malthusian growth rates $0.0849$, $0.0709$, $0.0644$
(Figure~\ref{fig:robust}B). Equal mean division time preserves neither
molecular recovery nor risk. Theorem~\ref{thm:direction} proves the sign for
all these clocks, but the quantitative certificate of
Proposition~\ref{prop:smooth} is for $m=1$.

\section{Large site counts: balanced founders under a smooth readout}
\label{sec:boundary}

Figure~\ref{fig:summary}A might suggest that the balanced-founder gap tends to
a positive limit. For a fixed smooth readout it does not. The reason is
that the symmetry line $a=r$ is invariant for the mean-field dynamics, a
founder near it stays near it for a time of order $\log N$, and on the line
the smooth hazard of our example makes the population subcritical.

Assume symmetric constants $w_A=w_R=w$, $k_A=k_R=k$, $e_A=e_R=e$,
$h_A=h_R=h$, and let the hazard be a function of the contrast,
\[
 d_N(a,r)=\varphi\bigl(c(a,r)\bigr),\qquad c(a,r)=\frac{a-r}N\in[-1,1],
\]
where $\varphi\ge0$ is Lipschitz with constant $\ell$ and satisfies
$\theta:=\varphi(0)-b>0$. Put
\[
 \psi=\varphi(0)-\inf\varphi,\qquad\Lambda=\max\{0,\,k-e-\tfrac34b\},\qquad
 \sigma^2=\max\{2w+k,\ e+\tfrac h2\}+\tfrac b2 .
\]

\begin{theorem}[Balanced founders die out as $N\to\infty$]\label{thm:boundary}
Under these assumptions, for either daughter law, every $N$, every founder
$x\in S_N$ and every $t\ge0$,
\begin{equation}\label{eq:fmbound}
 1-q_D(x)\le\E_xK(t)\le e^{-\theta t}\Bigl[1+\ell\,t\,e^{(\psi+\Lambda)t}\bigl(c(x)^2+\sigma^2t/N\bigr)^{1/2}\Bigr].
\end{equation}
Consequently, if $|c(x_N)|\le C_0N^{-1/2}$ and $\psi+\Lambda>0$, then with
$\gamma=\theta/[2(\psi+\Lambda)]$ and $t_N=\log N/[2(\psi+\Lambda)]$,
\begin{equation}\label{eq:rate}
 0\le q_I(x_N)-q_J(x_N)\le1-q_J(x_N)\le N^{-\gamma}\Bigl[1+\ell\,t_N\bigl(C_0^2+\sigma^2t_N\bigr)^{1/2}\Bigr]\xrightarrow[N\to\infty]{}0,
\end{equation}
where the first inequality requires $\varphi$ to be nonincreasing
\textup{(}Theorem~\ref{thm:direction}\textup{)} and the others do not.
\end{theorem}
\begin{proof}
Survival forever implies $K(t)\ge1$, so $1-q_D(x)\le\PP_x(K(t)\ge1)\le\E_xK(t)
=(e^{tA}\one)_x$, which by Lemma~\ref{lem:construction}(iv) is the same in the
two models. Write $A=\gen+\diag(b-d)$, where
\[
 \gen=Q+2b\,(D_1-I)
\]
is the generator of a conservative Markov chain $X$ on $S_N$: the molecular
chain together with a halving $x\mapsto Y$ at rate $2b$. By the Feynman--Kac
formula for finite chains (the many-to-one formula
\cite{bansaye2011,hardy2009}),
\[
 \E_xK(t)=\E_x\exp\Bigl(\int_0^t(b-d)(X_s)\,ds\Bigr)=e^{-\theta t}n_x(t),\qquad
 n(t)=e^{t(\gen+\diag\chi)}\one,\quad\chi=\varphi(0)-\varphi\circ c .
\]
Here $\chi\le\ell|c|$ and $\chi\le\psi$. Since $e^{t\gen}\one=\one$, the
Duhamel formula gives
\begin{equation}\label{eq:duhamel}
 n_x(t)=1+\int_0^t\E_x\Bigl[\chi(X_s)\exp\Bigl(\int_0^s\chi(X_\tau)\,d\tau\Bigr)\Bigr]ds
 \le1+\ell\int_0^te^{\psi s}\,\E_x|c(X_s)|\,ds .
\end{equation}
It remains to bound the second moment of the contrast along $X$. Under $Q$
the contrast jumps by $\pm1/N$. With symmetric constants the difference of
the drifts of $a$ and $r$ is
$u\,k(a-r)/N-e(a-r)$, the recruited-erasure terms $h\,ar/N$ cancelling, so
the drift of $c$ is $(k\,u/N-e)\,c$. The total jump rate is
\[
 u\bigl(2w+k\tfrac{a+r}N\bigr)+e(a+r)+2h\tfrac{ar}N\le u(2w+k)+(a+r)\bigl(e+\tfrac h2\bigr)\le N\max\{2w+k,\,e+\tfrac h2\},
\]
using $ar\le(a+r)N/4$. At a halving, $I-J$ has mean $(a-r)/2$ and variance
$(a+r)/4$, so $\E_xc(Y)^2=c^2/4+(a+r)/(4N^2)$. Therefore
\[
 \gen c^2=2\bigl(k\tfrac uN-e\bigr)c^2+\frac{\text{jump rate}}{N^2}+2b\Bigl(\frac{a+r}{4N^2}-\frac34c^2\Bigr)
 \le2\Lambda\,c^2+\frac{\sigma^2}N .
\]
By Dynkin's formula and Gronwall's inequality,
$\E_xc(X_s)^2\le e^{2\Lambda s}\bigl(c(x)^2+\sigma^2s/N\bigr)$. Insert
$\E|c|\le(\E c^2)^{1/2}$ into \eqref{eq:duhamel} and bound the integrand by
its value at $s=t$ to obtain \eqref{eq:fmbound}. For \eqref{eq:rate} take
$t=t_N$, so that $e^{(\psi+\Lambda)t_N}=N^{1/2}$ and
$e^{-\theta t_N}=N^{-\gamma}$.
\end{proof}

For the smooth readout \eqref{eq:smooth} with $s=8$ and the constants of
Table~\ref{tab:param} one has $\varphi(0)=0.155$, $\theta=0.055$,
$\ell=0.29s/4=0.58$, $\psi<0.145$, $\Lambda=0.915$, $\sigma^2=1.07$ and
$\gamma=0.055/2.12\approx0.0259$. The theorem covers the balanced founders
$(N/2,N/2)$, the square-root family
$(\lfloor(N+\sqrt N)/2\rfloor,\,N-\lfloor(N+\sqrt N)/2\rfloor)$, and any fixed
compact slope interval with the constants taken at the largest slope.

\begin{remark}[What the limit does and does not say]\label{rem:boundary}
(i) The bound \eqref{eq:rate} is asymptotic and the rate is slow; with the
constants above it is below one only for astronomically large $N$. This is in
the nature of the problem and not only of the proof: a founder near the line
must survive a subcritical phase whose length grows like $\log N$ before its
contrast is macroscopic, which suggests a small negative power of $N$. The
diagnostics agree: under the smooth readout the survival probability of
$(N/2,N/2)$ falls from $0.377$ at $N=4$ to $0.268$ at $N=128$
(Figure~\ref{fig:summary}B) and the balanced gap falls from $0.0156$ at
$N=8$ to $0.0087$ at $N=128$. A gap of order $10^{-2}$ at moderate $N$ and a
zero limit are therefore compatible, and the certificates of
Proposition~\ref{prop:smooth} are finite-size statements of independent
interest.

(ii) The argument uses only the common mean, so it cannot distinguish the two
models; it shows that both survival probabilities vanish.

(iii) Continuity of $\varphi$ at zero is essential. For the step readout
arbitrarily small contrasts of either sign change the hazard by $0.29$, the
bound $\chi\le\ell|c|$ fails, and in the diagnostics the balanced survival
probability increases with $N$ while the gap stays near $0.014$. The limit
for a discontinuous readout, or for a slope growing with $N$, requires a
fluctuation analysis that we do not attempt. Regularity of the biological
readout thus decides the asymptotic question.

(iv) If $\varphi(0)<b$ the population is supercritical on the symmetry line
and the theorem does not apply.
\end{remark}

\section{Substantial regrowth}\label{sec:regrowth}

Let $\tau_0$ be the extinction time and $\tau_L$ the first time the population
reaches $L\ge2$ cells, from one founder, and let $p_L(x)=\PP_x(\tau_L<\tau_0)$.

\begin{theorem}[Threshold identity]\label{thm:threshold}
Assume $d(x)\ge d_*>0$ on $S_N$. Then $\tau_0\wedge\tau_L<\infty$ almost
surely and
\begin{equation}\label{eq:identity}
 1-q_x=\E_x\Bigl[\Bigl(1-\prod_{j=1}^Lq_{X_j(\tau_L)}\Bigr)\one_{\{\tau_L<\tau_0\}}\Bigr],
\end{equation}
where $X_1(\tau_L),\dots,X_L(\tau_L)$ are the types present at $\tau_L$. If
$q_{\max}=\max_xq_x<1$,
\begin{equation}\label{eq:hitting}
 1-q_x\ \le\ p_L(x)\ \le\ \frac{1-q_x}{1-q_{\max}^{\,L}} .
\end{equation}
For $z\in(0,1)$ and a deadline $H>0$,
\begin{equation}\label{eq:deadline}
 \PP_x(\tau_L\le H)\ \ge\ 1-u_x(H;0)-z^{-(L-1)}\bigl[u_x(H;z)-u_x(H;0)\bigr].
\end{equation}
\end{theorem}
\begin{proof}
Before $\tau_0\wedge\tau_L$ the process lives on the finite set of
configurations with fewer than $L$ cells. From each of them, fewer than $L$
successive deaths within one time unit lead to extinction with a probability
bounded below uniformly, since $d\ge d_*$ and all other rates are bounded.
The strong Markov property gives a geometric tail for $\tau_0\wedge\tau_L$.
Divisions increase the size by exactly one, so the population at $\tau_L$ has
exactly $L$ cells, and given their types its extinction probability is the
product of their single-founder extinction probabilities. Extinction without
hitting $L$ has probability $1-p_L$, so
$q_x=1-p_L(x)+\E_x[\prod_jq_{X_j(\tau_L)};\tau_L<\tau_0]$, which is
\eqref{eq:identity}. Bounding the bracket between $1-q_{\max}^L$ and $1$ gives
\eqref{eq:hitting}. Finally $z^k\ge z^{L-1}$ for $1\le k<L$ gives
$\PP_x(1\le K(H)<L)\le z^{-(L-1)}[u_x(H;z)-u_x(H;0)]$; subtract this from
$1-\PP_x(K(H)=0)$ and use $\{K(H)\ge L\}\subseteq\{\tau_L\le H\}$.
\end{proof}
The inclusion at the end is not an equality: \eqref{eq:deadline} uses the
event $K(H)\ge L$ as a lower bound for first passage. After collecting terms
the coefficient of $u_x(H;0)$ is $z^{-(L-1)}-1>0$, so a certificate needs a
\emph{lower} enclosure of $u_x(H;0)$ and an \emph{upper} enclosure of
$u_x(H;z)$.

\begin{theorem}[Certified risk separations]\label{thm:application}
\textup{(a)} For the two-site step-readout source and one $AA$ founder,
\[
 p^J_{300}>0.779,\qquad p^I_{300}<0.681,\qquad\PP_J(\tau_{300}\le300)>0.765 ,
\]
so the probability of ever reaching $300$ cells differs by more than $0.098$
and the probability of reaching $300$ cells by time $300$ by more than
$0.084$.

\textup{(b)} For the $16$-site smooth source and one $(9,7)$ founder,
uniformly in $s\in[7.999,8.001]$,
\[
 p^J_{2000}>0.51438,\qquad p^I_{2000}<0.50138,\qquad p^J_{2000}-p^I_{2000}>0.013 ,
\]
and uniformly in $s\in[7,9]$, $p^J_{2000}-p^I_{2000}>0.0122$.
\end{theorem}
\begin{proof}
(a) Insert the tight boxes of Section~\ref{sec:instances},
$q_J(AA)<0.2207253$, $q_I(AA)>0.3351447$ and $\max_xq_I(x)<0.9875805$, into
\eqref{eq:hitting}: $p^J_{300}\ge1-q_J(AA)>0.779$ and
$p^I_{300}\le0.6648553/(1-0.9875805^{300})<0.681$. For the deadline take
$z=299/300$. Appendix~\ref{app:validation} proves
$u_{J,AA}(300;0)>0.2207177$ and $u_{J,AA}(300;299/300)<0.2259640$, and
\eqref{eq:deadline} gives
\[
 \PP_J(\tau_{300}\le300)\ge1-0.2207177-(300/299)^{299}\,(0.2259640-0.2207177)>0.765 .
\]
The eventual bound
$p^I_{300}$ also bounds $\PP_I(\tau_{300}\le300)$. With the coarse constants
$q_J(AA)\le0.235$, $q_I(AA)\ge0.313291$, $q_{I,\max}<0.988$,
$u_J(300;0)>0.22$ and $u_J(300;z)<0.227$ the same computation gives the weaker
values $0.765$, $0.706$ and $0.760$; these are the versions whose arithmetic
is checked in Lean.

(b) Use Proposition~\ref{prop:smooth}: in case (a) of that proposition,
$q_J(9,7)<0.48562$, $q_I(9,7)>0.49988$, $q_{I,\max}<0.997$ and
$0.50012/(1-0.997^{2000})<0.50138$. For $[7,9]$ the same computation is
carried out exactly on each of the $40$ slope cells with the boxes of that
cell, and the smallest resulting gap exceeds $0.0122$.
\end{proof}

These are unconditional probabilities for one living founder of the declared
type at time zero. There is no preceding treatment phase and no selection on
later survival; the $16$-site statement is eventual and no finite deadline is
asserted for it. Thresholds of hundreds or thousands of cells establish a
substantial clonal-count event within a dilute branching model; they do not
calibrate a relapse and do not validate independent growth at such
densities.

\section{What observations can identify}\label{sec:observation}

\subsection{Paired sister clones}
Enrol sister pairs prospectively at birth from a specified mother
preparation, culture the two daughters separately, and record whether each
clone is extinct at follow-up time $H$. Given the daughter states, the two
indicators are independent with probabilities $f_H(Y),f_H(Z)$, where
$f_H(x)=u_x(H;0)$. Hence, for a fixed complete mother state $x$,
\[
 \Cov\bigl(\text{two extinction indicators}\mid X=x\bigr)=\Cov_x\bigl(f_H(Y),f_H(Z)\bigr)\le0
\]
under complementary allocation, by Lemmas~\ref{lem:embedding} and
\ref{lem:assoc}, and it is zero under the independent comparator. This
measures continuation fate and not a thresholded molecular reporter.

Pooling over heterogeneous mothers can reverse the sign. For an imperfectly
measured parental stratum $W$, total covariance gives
\[
 C_W=\E[c_X\mid W]+\Var(\mu_X\mid W),\qquad\mu_X=D_1f_H(X),\quad c_X=\Cov_X(f_H(Y),f_H(Z)).
\]
For example, for newborn ``active'' indicators $\one_{\{a>r\}}$ in the
two-site model, an $AA$ mother has covariance $1/2-(3/4)^2=-1/16$ and an $RR$
mother has covariance zero, but an equal mixture of $AA$ and $RR$ mothers has
pooled covariance $1/4-(3/8)^2=+7/64$. If $\mu_X$ has range at most $\Delta$
within the stratum, then $\Var(\mu_X\mid W)\le\Delta^2/4$. If each binary fate
is misclassified with probability at most $\zeta$, the joint probability and
the product of marginals each change by at most $2\zeta$. If a fraction at
most $\epsilon$ of \emph{prospectively enrolled} pairs is missing, each of the
three probabilities changes by at most $\epsilon$. The average within-mother
covariance therefore lies in
\begin{equation}\label{eq:observation}
 \bigl[\,C_{\mathrm{obs}}-\Delta^2/4-4\zeta-3\epsilon,\ \ C_{\mathrm{obs}}+4\zeta+3\epsilon\,\bigr].
\end{equation}
This is a deterministic sensitivity interval. A sampling radius is separate:
for $n$ independent enrolled pairs, Hoeffding's inequality
\cite{hoeffding1963} applied to the two marginal frequencies and the joint
frequency, with a union bound, puts the empirical covariance within
$3\sqrt{\log(6/\alpha)/(2n)}$ of its expectation with probability at least
$1-\alpha$. The sampling unit is the independent pair and not the descendant
cell. Without justified bounds on $\Delta$ and on missingness a pooled
covariance does not identify partition as its cause; unrecorded enrolment
cannot be repaired afterwards, and sampling extant cells rather than
divisions introduces a growth bias that this analysis does not cover.

\subsection{Clone counts under independent capture}
High-throughput lineage barcoding
\cite{gutierrez2021,kong2020,biddy2018,oren2021} reports clone sizes after
incomplete sampling and usually does not identify the two immediate sister
branches. Suppose that a clone has $K(t)$ living cells and each is captured
independently with a known probability $p$, giving a count
$C\mid K\sim\Bin(K,p)$.

\begin{corollary}[Thinned clone counts]\label{cor:capture}
For either model,
$\E_xz^C=u_x(t;1-p+pz)$, $\E_xC=p\,m_x(t)$ and
$\Var_xC=p^2\Var_xK(t)+p(1-p)m_x(t)$. Under the hypotheses of
Theorem~\ref{thm:direction}, for a common founder law and a common $p$,
\[
 \PP(C_J=0)\le\PP(C_I=0),\qquad\E C_J=\E C_I,\qquad\Var C_J\le\Var C_I .
\]
\end{corollary}
\begin{proof}
$\E[z^C\mid K]=(1-p+pz)^K$; then apply Theorem~\ref{thm:direction} at
$z'=1-p\in[0,1]$ and Corollary~\ref{cor:variance}.
\end{proof}
Nondetection is not extinction: $\PP(C=0)=u(t;1-p)\ge u(t;0)$. Independent
thinning is an explicit assumption. Sequencing to a fixed total depth induces
dependence between clones, and state-dependent capture, barcode collisions or
silencing and the passage from reads to cells need their own observation
model. Clone-size dispersion is affected by many mechanisms, so
Corollary~\ref{cor:capture} is a prediction for matched founder preparations
and mean growth; it does not make barcode dispersion a substitute for the
paired-clone covariance, which is the statistic tied directly to the
mechanism and which requires the mother preparation and the two daughter
branches to be identified, for example by imaging, indexed isolation or
sufficiently resolved lineage recording.

\section{Scope, formalization and reproducibility}\label{sec:scope}

\paragraph{Chain of results.} Molecular order and conserved allocation fix
the sign (Theorem~\ref{thm:direction}), also for the variance
(Corollary~\ref{cor:variance}) and strictly so under irreducibility
(Proposition~\ref{prop:strict}). Continuation-risk covariance and a stable
population response fix the magnitude (Theorem~\ref{thm:response}). Exact
boxes give finite-size contrasts between preparations
(Proposition~\ref{prop:smooth}) and consequences at large population
thresholds (Theorem~\ref{thm:application}), while
Theorem~\ref{thm:boundary} shows that, for a smooth readout, the balanced
contrast is a finite-size effect. The magnitude is not universal: it depends
strongly on preparation, on the regularity of the readout and on the division
clock.

\paragraph{What is not claimed.} There is no unconditional advantage of
conservation: the theorem needs the stated order, a hazard monotone in it and
a division clock independent of the molecular state, and
Section~\ref{sec:statedep} shows where the proof breaks otherwise. Nonmonotone
hazards and fitness trade-offs are outside the theorem. Changing a symmetric
erasure rate removes marks of both kinds, so no monotonicity in that
parameter follows. The source supplies no empirical locus size, drug hazard
or competition law; hazards, site counts and clocks are exploratory, and no
clinical calibration is implied.

\paragraph{Formalization boundary.} \sloppy A Lean~4/mathlib development
(\texttt{PhenotypeMemory}, fourteen modules, compiled with warnings as errors
and without \texttt{sorry}; its only axioms are \texttt{propext},
\texttt{Classical.choice} and \texttt{Quot.sound}) checks: normalization of
the partition weights for all $(a,r)$; the two-site source formulas; the
supersolution inequalities for $\bar q_I$ and, uniformly for
$e\in[0.0099,0.0101]$, for $\bar q_J$; the bound $\bar q_I<0.988$; the eight
downward iterations and $q_I(AA)>0.31$; the weighted Jacobian inequality with
$\kappa=0.796$; reusable algebra for the state order, the one-bit
conditional-covariance step of Lemma~\ref{lem:assoc}, the response identity
and the threshold inequalities; and the coarse arithmetic
$0.686709/(1-0.988^{300})<0.706$ and
$1-0.22-(300/299)^{299}(0.227-0.22)>0.760$. The root declaration is
\begin{center}\small\texttt{PhenotypeMemory.continuation\_finite\_certificate}.\end{center}
Everything else is conventional: the chronological construction, the
couplings, the induction over allocation bits, the comparison theorem for
the generating functions, Corollary~\ref{cor:variance},
Propositions~\ref{prop:strict} and \ref{prop:obstruction},
Theorem~\ref{thm:boundary}, the strong Markov argument, the soundness of the
$153$-state rational checker and of the validated integrator, and the tight
two-site boxes. Exact computation is not presented as an end-to-end formal
probability theorem.

\paragraph{Reproducibility.} The source archive contains the scripts that
produce every number: \path{slope_certificate.py} (Proposition~\ref{prop:smooth};
about two seconds per slope cell), \path{deadline_validated.py}
(Appendix~\ref{app:validation}; about eighty seconds),
\path{check_paper.py} (an independent exact replay of all two-site constants,
the threshold arithmetic, the constants of Theorem~\ref{thm:boundary} and the
small exact examples), \path{diagnostics.py} and \path{make_figures.py}
(floating-point figures). Rational boxes are stored as exact fractions, and
the printed decimals are checked against them.

\appendix
\section{Validated microscopic generating functions}\label{app:validation}

Write the six-state joint vector field \eqref{eq:ODE} as
$\Phi(v)=c_0+\mathsf Lv+\mathsf P(v,v)$ with rational coefficients, where
$c_0=d$, $\mathsf L=Q-\diag(d)-bI$ and $\mathsf P=bD_J$. Here
$\|Q\|_\infty=1.06$, $d\le0.3$, $b=0.1$, so in the complex ball
$\|v\|_\infty\le2$,
\[
 \|\Phi(v)\|_\infty\le1.06\cdot2+0.3\,(1+2)+0.1\,(4+2)=3.62<4 .
\]
\emph{Local analyticity.} Let $v_0\in[0,1]^6$. The ball of radius one about
$v_0$ lies in $\{\|v\|_\infty\le2\}$. By the Cauchy--Picard theorem for
analytic vector fields the solution through $v_0$ extends analytically to
complex times as long as it stays in that ball, and along any ray it cannot
leave the ball before time $1/3.62>1/4$. Thus it is analytic and bounded by
$2$ on $|t|\le1/4$, Cauchy's estimate bounds its $j$th Taylor coefficient by
$2\cdot4^j$, and with step $\Delta=1/50$ the truncation after degree $20$ has
error at most
\[
 \eta=\frac{2\,(2/25)^{21}}{1-2/25}<2.01\times10^{-23}.
\]
\emph{Arithmetic.} The Taylor coefficients obey the exact recurrence
$(j+1)c_{j+1}=\mathsf Lc_j+\one_{\{j=0\}}c_0+\sum_{i=0}^j\mathsf P(c_i,c_{j-i})$.
Every operation is enclosed by integer intervals at denominator $10^{40}$,
the polynomial is evaluated by an interval Horner scheme, and the new centre
is the rounded midpoint projected coordinatewise to $[0,1]$. Projection does
not increase the distance to the exact solution, which lies in the cube. The
largest distance from the new centre to an endpoint of the enclosure, over
all $15{,}000$ steps and both solves, is
$\varepsilon_{\mathrm{rd}}=10^{-40}$ (one unit of the last place); this number
is recorded by the code.

\emph{Propagation.} On the real cube the Jacobian of $\Phi$ has nonnegative
off-diagonal entries and row sums $-d_x-b+b\sum_y\partial_yD_J(v,v)_x\le-d_x+b\le0.09$,
so its logarithmic $\infty$-norm is at most $0.09$ \cite{soderlind2006}, and
two solutions that remain in the cube, which is forward invariant, separate
by at most a factor $e^{0.09t}$. The usual telescoping argument for one-step
methods \cite{nedialkov1999} bounds the global error at $t=300$ by
\[
 e^{27}\bigl\{15000\,(\eta+\varepsilon_{\mathrm{rd}})+10^{-40}\bigr\}<1.8046\times10^{-7},
\]
where $e^{27}<6\times10^{11}$ is proved by a rational degree-$200$ Taylor sum
with a geometric tail. Almost all of the budget is the analytic truncation
term; the arithmetic term contributes less than $10^{-23}$.

\begin{center}
\begin{tabular}{ccc}
\toprule
initial value $z$ & lower bound for $u_{J,AA}(300;z)$ & upper bound for $u_{J,AA}(300;z)$\\\midrule
$0$ & $0.2207177851$ & $0.2207181461$\\
$299/300$ & $0.2259636331$ & $0.2259639941$\\\bottomrule
\end{tabular}
\end{center}
The integrator's coefficients are checked independently against the rational
source and partition formulas. No floating-point solver tolerance enters an
enclosure. The enclosure for $z=0$ is consistent with
$u_{J,AA}(300;0)\le q_J(AA)<0.2207253$.

\section{Certificate data and proof dependencies}\label{app:data}

The files \path{data/slope_certificate_7.999_8.001_1.json} and
\path{data/slope_certificate_7_9_40.json} contain, per slope cell, the rational
risk boxes at $(9,7)$, the upper bounds on the all-active gap and on
$\max q_I$, iteration counts and smallest supersolution residuals; the
single-cell file also contains all $153$ hazard brackets and both full
$153$-coordinate boxes. \path{data/deadline_validated.json} holds the
enclosures of Appendix~\ref{app:validation} as fractions,
\path{data/paper_checks.json} the replayed two-site constants, and
\path{data/diagnostics.json} the floating-point diagnostics behind the
figures.
\begin{center}\small
\begin{tabular}{p{.27\linewidth}p{.65\linewidth}}
\toprule
Conclusion & Ingredients\\\midrule
Direction, all $N$ & Monotone molecular coupling; ordered allocation; Harris inequality on reversed bits; monotone iteration; Kamke comparison.\\
Variance ordering & Second factorial moment equation; increasing mean; same association inequality; nonnegative semigroup.\\
Strictness & Irreducibility; classical extinction criterion; one-bit covariance bound; response identity.\\
Magnitude & Resolvent; covariance forcing; difference of squares; weighted contraction; rational boxes.\\
Smooth $N=16$ contrast & Rational exponential bounds; statewise slope envelopes; hazard coupling; lower iteration; supersolutions.\\
Large-$N$ limit & Common mean; Feynman--Kac; Duhamel; second moment of the contrast; Gronwall.\\
Regrowth & Almost-sure stopping; strong Markov; extinction boxes; validated generating functions for the deadline.\\
Observation & Conditional independence; total covariance; range--variance bound; Hoeffding; binomial thinning.\\\bottomrule
\end{tabular}
\end{center}

\bibliographystyle{abbrv}
\bibliography{refs}
\end{document}
